\documentclass[12pt]{article}

\input{../../_assets/preambulo.tex}


\begin{document}

    % 1. Foto de fondo
    % 2. Título
    % 3. Encabezado Izquierdo
    % 4. Color de fondo
    % 5. Coord x del titulo
    % 6. Coord y del titulo
    % 7. Fecha

    
    \input{../../_assets/portada}
    \portadaExamen{ffccA4.jpg}{Ecuaciones\\Diferenciales II\\Examen III}{Ecuaciones Diferenciales II. Examen III}{MidnightBlue}{-8}{28}{2026}{David Muñoz Gómez}

    \begin{description}
        \item[Asignatura] Ecuaciones Diferenciales II.
        \item[Curso Académico] 2024/25.
        \item[Grado] Doble Grado en Ingeniería Informática y Matemáticas.
        \item[Grupo] Único.
        \item[Profesor] Rafael Ortega Ríos.
        \item[Descripción] Convocatoria extraordinaria.
        \item[Fecha] 11 de julio de 2025.
        % \item[Duración] ---.
    
    \end{description}
    \newpage


    % ------------------------------------
    
    \begin{ejercicio}
        Se supone dada una sucesión de funciones $f_n : [0, 1] \to \bb{R}$ que converge uniformemente en $[0, 1]$. Además cada función $f_n$ es continua. ¿Se puede asegurar que la sucesión $\{f_n\}$ es equicontinua?
    \end{ejercicio}

    \begin{ejercicio}
        Demuestra que la solución maximal del problema de Cauchy
        \[ \dot{x} = \frac{1}{x-1} + e^t, \quad x(0) = 2 \]
        cumple $\omega = +\infty$.
    \end{ejercicio}

    \begin{ejercicio}
        El campo $X : \bb{R} \times \bb{R}^d \to \bb{R}^d$ es continuo y se supone que las funciones $\phi : ]-\infty, 1] \to \bb{R}^d$, $\varphi : [1, +\infty[ \to \bb{R}^d$ son soluciones de la ecuación $\dot{x} = X(t, x)$. Además $\phi(1) = \varphi(1)$, $\phi(0) = 0$. Prueba que el problema
        \[ \dot{x} = X(t, x), \quad x(0) = 0 \]
        admite una solución maximal definida en toda la recta real.
    \end{ejercicio}

    \begin{ejercicio}
        Se considera la ecuación
        \[ \dot{x} = 2x - 1. \]
        Encuentra sus puntos de equilibrio y discute sus propiedades de estabilidad.
    \end{ejercicio}

    \begin{ejercicio}
        Para cada $\veps > 0$ se considera la solución maximal $x_{\veps}(t)$ del problema
        \[ \veps \dot{x} = \operatorname{sen}(\veps^2 t)(x - 1), \quad x(0) = 2. \]
        Calcula, si existe, $\lim\limits_{\veps \to 0^+} x_{\veps}(32)$.
    \end{ejercicio}

    \begin{ejercicio}
        Demuestra que la ecuación integral
        \[ x(t) = 7 + \frac{1}{4} \int_{t-1}^{t+1} \cos x(s) \, ds \]
        admite a lo sumo una solución continua y acotada $x : \bb{R} \to \bb{R}$.
    \end{ejercicio}

    \newpage
    \setcounter{ejercicio}{0} % Reiniciar contador de ejercicios
    % \noindent
    % \textbf{Solución.}

\end{document}